\RequirePackage{pdfmanagement}
\SetKeys[document/metadata]{pdfversion=1.7, pdfstandard=A-2u, lang=en-GB}
\input{preamble/pdfA-compliance/glyphtounicode}
\documentclass[11pt,a4paper]{report}

\newcommand*{\PaperTitle}{Criticality of the Hopf bifurcations in the Becker--Döring equations with atomization}
\newcommand*{\PaperTitlePlain}{Criticality of the Hopf bifurcations in the Becker-Doring equations with atomization}
\newcommand*{\PaperShortTitle}{Hopf criticality in Becker--Döring with atomization}
\newcommand*{\PaperSubtitle}{chemistry-open-problems \textperiodcentered{} 18 October 2026 \textperiodcentered{} partial result, preliminary}
\newcommand*{\PaperCredits}{Solver: Claude Opus 5.5 (effort xhigh), sub-agent of the orchestrator; citing-works search by a nested Claude Opus 5.5 sub-agent (effort xhigh) \textperiodcentered{} LaTeX edition: Claude Opus 5.5 (effort xhigh) \textperiodcentered{} Independent verification: pending}
\newcommand*{\PaperKeywords}{Becker--Döring equations, nucleation, Hopf bifurcation, first Lyapunov coefficient, interval arithmetic}
\input{chemop-preamble}

\begin{document}
\makepapertitle

\begin{paperabstract}
Pego and Velázquez proved that the Becker--Döring equations with atomization undergo Andronov--Hopf bifurcations
at \(\kappa=\kappa_j(M)\) for large maximal cluster size \(M\), and wrote that they had not managed to determine
analytically whether the bifurcating periodic solutions are stable. We evaluate the first Lyapunov coefficient
\(\ell_1\) in exact closed form at a cost independent of \(M\), enclose the Hopf points by a Krawczyk operator in ball
arithmetic, and certify the sign of \(\ell_1\) for every integer \(M\) in \(15\le M\le10^5\) on the first branch
(\(44\le M\le10^5\) and \(87\le M\le10^5\) on the second and third): the bifurcation is supercritical, with stable
periodic orbits, except at \(M=15\) on the first branch, where it is subcritical. A formal reduction for
\(M\to\infty\) to a scalar renewal equation predicts a negative Landau coefficient on every branch, and certified values up
to \(M=10^8\) approach the predicted limit. The infinite range \(M>10^5\) remains unproved.
\end{paperabstract}

\section{Problem and source}

R. L. Pego, J. J. L. Velázquez, \emph{Temporal oscillations in Becker--Döring equations with atomization},
Nonlinearity 33 (2020) 1812--1846, doi:10.1088/1361-6544/ab6815, arXiv:1905.02605 (v2 read in full, LaTeX source).

Theorem 3.3 of the paper states that for \(k\in\mathbb N\) and \(M>M_{0,k}\) the system admits a Hopf bifurcation
as \(\kappa=K\sqrt M\) passes through each \(\kappa_j(M)\), \(j=1,\dots,k\). Immediately after it the authors write
(verbatim, §3.3):

\begin{quote}
``We have not managed to determine analytically whether the bifurcating solutions are stable (the supercritical
case) or not. Many of our numerical computations, as in Fig. 1, are consistent with the presence of stable
periodic solutions, however.''
\end{quote}

The constant \(M_{0,k}\) is not explicit.

Novelty check (2026-10-18): 20 citing works (OpenAlex and Semantic Scholar) were checked; every arXiv full text was
grepped. Later work by the authors and searches on key phrases were also checked. Nothing settles the question.
The closest work is Niethammer--Pego--Schlichting--Velázquez (SIAM J. Appl. Math. 2022, arXiv:2102.06751). It
carries out a \emph{formal} direction-of-bifurcation computation for a \emph{different} model (Becker--Döring with injection
and depletion). Details are in \texttt{references.md}.

\section{Model, linearization, exact quadratic structure}

The unknowns are \(n_\ell\ge0\), \(1\le\ell\le M\), with \(N=M-1\) and \(J_\ell=n_\ell n_1-n_{\ell+1}\) for
\(1\le\ell\le N\) (paper (1)--(4)):
\[\dot n_\ell=J_{\ell-1}-J_\ell\ (2\le\ell\le N),\qquad \dot n_M=J_N-Kn_M,\qquad
\dot n_1=-J_1-\sum_{\ell=1}^N J_\ell+MKn_M .\]
The mass \(\sum_\ell \ell n_\ell\) is conserved.

The constant state \(n_\ell\equiv A:=1+K\) is an equilibrium. Write \(n=A\mathbb 1+v\). A direct expansion
(checked in \texttt{scripts/01\_check\_model.py}, claims M1--M3) gives the \textbf{exact} identity
\[\dot v = Bv + v_1\,Sv ,\tag{2.1}\]
with no cubic or higher terms. Here \(B\) is the paper's matrix (27). The linear map \(S\) is
\[(Sv)_1=-2v_1-\sum_{\ell=2}^{N}v_\ell,\qquad (Sv)_\ell=v_{\ell-1}-v_\ell\ (2\le\ell\le N),\qquad (Sv)_M=v_N .\]
Equivalently, \(f(n)=(n_1-1)\,Sn+Tn\) with \(T\) linear. The nonlinearity is the product of the ``transport speed''
\(c=n_1-1\) and a fixed linear operator.

Mass conservation: \(\boldsymbol\ell=(1,2,\dots,M)\) satisfies \(\boldsymbol\ell B=0\) and \(\boldsymbol\ell S=0\).
Hence the hyperplane \(H=\{\boldsymbol\ell\cdot v=0\}\) is invariant under (2.1), and \(\mathrm{Ran}\,B=H\).

By paper Lemma 3.1, \(0\) is a simple eigenvalue of \(B\), with eigenvector \(\bar v\notin H\). Hence \(B|_H:H\to H\)
is invertible.

Kuznetsov's multilinear forms for (2.1) are
\[\mathcal B(x,y)=x_1Sy+y_1Sx,\qquad \mathcal C\equiv0 .\]

\section{Hopf setting and first Lyapunov coefficient}

\textbf{Setting.} As in the paper, the bifurcation parameter is \(\kappa\), and the dynamics is restricted to the
invariant affine hyperplane through \(A\mathbb 1\). In the coordinate \(v\) this is \(H\), a fixed space, and the
equilibrium is \(v=0\) for every \(\kappa\).

\textbf{Eigenvalues.} The eigenvalues of \(B|_H\) are those of \(B\) except the simple eigenvalue 0. At a Hopf point
\(\kappa_j\), \(B\) has eigenvalues \(\pm i\omega\) with \(\omega>0\), and \(\mathrm{Re}\,d\lambda/d\kappa>0\)
(paper (33), re-certified in §5).

\textbf{Coefficient.} Let \(Bq=i\omega q\), let \(u\) be a row vector with \(uB=i\omega u\), normalized by
\(uq=1\), and set \(\langle p,x\rangle:=ux\). Kuznetsov's formula (\emph{Elements of Applied Bifurcation Theory}, 3rd
ed., §3.5), with \(\mathcal C=0\), reads
\[\ell_1=\frac1{2\omega}\,\mathrm{Re}\Big[-2\,u\,\mathcal B\big(q,w_{11}\big)+u\,\mathcal B\big(\bar q,w_{20}\big)\Big],\]
\[w_{11}=(B|_H)^{-1}\mathcal B(q,\bar q)\in H,\qquad w_{20}=(2i\omega-B)^{-1}\mathcal B(q,q). \tag{3.1}\]

\textbf{Why the full-space objects may be used.}
- \(q\in H\), because \(\boldsymbol\ell q=\boldsymbol\ell Bq/(i\omega)=0\).
- \(\mathcal B(x,y)\in H\) for all \(x,y\), because \(\boldsymbol\ell S=0\).
- \(u|_H\) is the adjoint eigenvector of \(B|_H\).
- \(w_{20}\in H\) automatically, because \(H\) is invariant under \(B\) and \(2i\omega\notin\mathrm{spec}\,B\).
- \(w_{11}\) must be taken as the unique solution of \(Bw=\mathcal B(q,\bar q)\) \textbf{in \(H\)}. Adding a multiple of
\(\bar v\) would change \(\ell_1\).

Formula (3.1) is basis-free (projection method), so it is the first Lyapunov coefficient of the restricted system.

\textbf{Sign convention.} \(\ell_1\) depends on the normalization of \(q\) only through a positive factor. Two
normalizations are used: \(q^Hq=1\) (``unit'') and \(q_1=1\).

Since \(\mathrm{Re}\,\lambda'(\kappa_j)>0\):
- \(\ell_1<0\) means \textbf{supercritical}: periodic orbits exist for \(\kappa>\kappa_j\), \(\kappa\approx\kappa_j\),
with amplitude \(\propto(\kappa-\kappa_j)^{1/2}\), and they attract within the centre manifold.
- \(\ell_1>0\) means \textbf{subcritical}.

For \(j=1\) a supercritical orbit is orbitally asymptotically stable within \(H\). For \(j\ge2\) the equilibrium
already has \(j-1\) unstable pairs, so the bifurcating orbits are unstable in the full system whatever the sign.
The criticality is then a centre-manifold statement.

\textbf{Remark (intrinsic).} \(\ell_1\) depends only on the vector field at the critical point. The conclusion therefore
holds for any one-parameter path through the Hopf point in the \((K,\text{mass})\) plane that crosses transversally.
An example is \(K\) varying at fixed total mass, where the equilibria are the non-constant ones of paper (15).

\section{\texorpdfstring{Exact closed-form evaluation of \(\ell_1\) for any \(M\)}{Exact closed-form evaluation of \textbackslash ell\_1 for any M}}

All vectors occurring in (3.1) have the form
\[v_\ell=\sum_{e} c_e\,\beta_e^{\ell}+\sum_{k\in P}d_k\,\delta_{\ell k},\qquad P\subset\{1,2,N,M\},\tag{4.1}\]
with finitely many bases \(\beta_e\). Each base is a monomial \(\varphi^{e_0}\bar\varphi^{e_1}A^{e_2}s^{e_3}\bar s^{e_4}\)
with integer exponents; \(\varphi\) and \(s\) are defined below. All operations on such vectors are exact and cost
\(O(1)\) in \(M\). Implementation: \texttt{scripts/bdlib.py}, classes \texttt{Problem} and \texttt{Seq}.

\textbf{Right eigenvector} (paper §5). For \(\lambda=(A-\varphi^{-1})(\varphi-1)\) and \(\psi=1/(A\varphi)\):
\[q_\ell=c_1\varphi^{M-\ell}+c_2\psi^{M-\ell},\qquad (c_1,c_2)=(f(\psi),-f(\varphi)),\qquad f(\varphi)=A\varphi^{M-1}+1-\varphi^{-1},\]
valid for \(1\le\ell\le M\) whenever \(\varphi\) is a root of the paper's \(F\) (eqs. (71), (76)).

\textbf{Left eigenvector.} The transposed equations, columns \(2\le k\le N\) of \(uB=\lambda u\), read
\[u_{k-1}-(A+1+\lambda)u_k+Au_{k+1}=Ku_1 .\]
Hence \(u_k=a_1\varphi^k+a_2\psi^k+u_c\), with \(u_c=-Ku_1/\lambda\) and
\(u_1=(a_1\varphi+a_2\psi)\lambda/(\lambda+K)\). Columns 1 and \(M\) give a singular \(2\times2\) system for
\((a_1,a_2)\); one row of it fixes the null vector.

\textbf{Linear maps.} \(B\) and \(S\) act on a term \(\beta^\ell\) as multiplication by \(A/\beta-(A+1)+\beta\) (bulk rows
of \(B\)) or \(1/\beta-1\) (bulk rows of \(S\)). Rows \(1\) and \(M\), and the neighbours of point corrections, are
evaluated exactly. The full sums in row 1 are geometric sums. The differences are stored as point corrections.

\textbf{Solving \((B-\mu)w=r\)} for \(r\) of the form (4.1) (here \(\mu=0\) or \(2i\omega\)). Ansatz:
\[w=\sum_e \frac{c_e}{A/\beta_e-(A+1)+\beta_e-\mu}\,\beta_e^\ell+h_1s_1^\ell+h_2s_2^\ell+e_1\delta_{\ell1}+e_M\delta_{\ell M},\]
where \(s_{1,2}\) are the roots of \(s^2-(A+1+\mu)s+A=0\):
- for \(\mu=0\): \(s\in\{1,A\}\);
- for \(\mu=2i\omega\): one new generator \(s\), the other root being \(A/s\).

The bulk rows \(3\le\ell\le N-1\) hold identically. Rows \(2,N\) (which involve \(w_1,w_M\)) and rows \(1,M\) give
four linear equations for \((h_1,h_2,e_1,e_M)\).

\emph{Why this family contains the true solution.} Rows \(2,\dots,N\) form a three-term recurrence. Their solution set
is an affine space of dimension 2, and it is the 2-dimensional subfamily of the ansatz fixed by rows \(2\) and
\(N\).

\emph{The case \(\mu=0\).} Row 1 is redundant, since \(\boldsymbol\ell\,(Bw-r)=0\) whenever rows \(2,\dots,M\) hold. It
is replaced by the condition \(\sum_\ell \ell w_\ell=0\), which selects \(w_{11}\in H\). The \(4\times4\) system is
solved in ball arithmetic (\texttt{acb\_mat.solve}). Success proves that it is nonsingular, so the computed ball contains
the exact solution. The resulting point corrections are numerically zero (\(\sim10^{-35}\)), as they should be when
\(r\) has no correction at rows 2 and \(N\).

\textbf{Inner products.} \(\sum_{\ell=1}^M u_\ell v_\ell\) is computed with exact geometric sums
\(\sum\beta^\ell=(\beta^{M+1}-\beta)/(\beta-1)\). A base that is identically 1 (zero exponent vector, e.g.
\(\psi\cdot A\varphi\)) contributes the exact count \(M\). Products \(\beta^\ell\) of size \(A^{\pm M}\) are harmless
in Arb's arbitrary-exponent balls.

\textbf{Cross-checks.}
- The closed form agrees with a dense floating-point implementation of (3.1) to \(10^{-12}\)--\(10^{-14}\) relative
error for \(M=20,\dots,200\), \(j\le4\) (\texttt{01\_check\_model.py}, M6).
- The eigenvector residuals are at rounding level.

\section{Certification}

\textbf{Hopf points.} Unknowns \(\kappa=K\sqrt M\) and \(z=M(\varphi-1)\) (the paper's scaling). The system is
\[Q^\varepsilon(z;\kappa):=K^{-3}F(1+z/M)=0\ (\text{complex}),\qquad \mathrm{Re}\,\Lambda=0,\quad \Lambda=M\lambda/K .\]
This is 3 real equations in 3 real unknowns, with \(F=f_1g_2-g_1f_2\) exactly as in the paper, eqs. (71)--(76).
- Derivatives come from forward-mode dual numbers over Arb balls.
- The Krawczyk operator \(\mathcal K(X)=m-Yf(m)+(I-YJ(X))(X-m)\) is evaluated on a box \(X\) of relative radius
\(10^{-25}\). The inclusion \(\mathcal K(X)\subset\mathrm{int}\,X\) proves existence and uniqueness of a zero in
\(X\).
- By paper Prop. 5.3 a zero is an eigenvalue \(\lambda=(A-\varphi^{-1})(\varphi-1)\) of \(B\), which is purely
imaginary at the zero (\(\mathrm{Re}\,\Lambda=0\)). It is not spurious, since \(\mathrm{Im}\,\varphi\neq0\).

\textbf{Nondegeneracy (certified at every point).}
- \emph{Simplicity:} \(0\notin\partial_zQ^\varepsilon\) on the box. Then \(\varphi\) is a simple root of \(F\), and by
paper Lemma 5.4 \(i\omega\) is a simple eigenvalue of \(B\).
- \emph{Transversality:} \(\mathrm{Re}\big(\Lambda+\kappa\,d\Lambda/d\kappa\big)>0\) along the root curve, with
\(dz/d\kappa=-\partial_\kappa Q^\varepsilon/\partial_zQ^\varepsilon\). This equals \(M^{3/2}\,\mathrm{Re}\,d\lambda/d\kappa\).

\textbf{Sign of \(\ell_1\).} The closed form of §4 is evaluated with \(K=\kappa/\sqrt M\) and \(\varphi=1+z/M\) given as
the Krawczyk enclosure. The formulas are exact identities at the true Hopf point, so the output ball contains the
true \(\ell_1\). Precision is 192 bits. The Krawczyk boxes have relative radius \(10^{-25}\), and the resulting
\(\ell_1\) balls relative radius \(10^{-36}\)--\(10^{-29}\) (measured at \(M=15\) and \(M=10^5\)). The sign is
therefore never in doubt.

\textbf{Remaining hypothesis H2} (no other eigenvalue of \(B|_H\) on \(i\mathbb R\)).
- For \(M\le120\) it is certified with Arb's rigorous eigensolver (\texttt{acb\_mat.eig}, \texttt{07\_spectrum\_H2.py}). All \(M\)
eigenvalues of \(B(K)\), with \(K\) the certified ball, are enclosed and proved simple. Exactly one enclosure
contains 0, exactly two meet \(i\mathbb R\) (at \(\pm i\omega\)), \(2(j-1)\) have positive real part and the rest
negative real part.
- This also proves rigorously that the point is the \(j\)-th crossing.
- Result: all 217 points with \(M\le120\) PASS (106 for \(j=1\), 77 for \(j=2\), 34 for \(j=3\));
log in \texttt{certificates/logs/07\_spectrum\_H2.log}.
- For larger \(M\), H2 and the branch label \(j\) are supported by floating-point eigenvalue counts: at
\(\kappa_j(1\mp10^{-7})\), \(j-1\) resp. \(j\) unstable pairs, for all \(M\le300\) and every 50th \(M\le2000\).
The paper's Theorem 3.2 gives H2 for \(M>M_{0,k}\).

\textbf{Generation and independent re-check.}
- The sweep files are generated by \texttt{03\_certify\_range.py} (continuation in \(M\)).
- \texttt{04\_verify\_certificates.py} then re-proves every line from the stored midpoints alone, without continuation:
Krawczyk, \(\ell_1\) enclosure, simplicity, transversality. It also checks that every integer \(M\) of the range is
present.
- Done for every line of all six files (\(3\times(\approx10^4+9\cdot10^4)\) Hopf points). Logs are in
\texttt{certificates/logs/04\_verify\_*.log}.

\section{Results on the finite range}

\textbf{Theorem A (computer-assisted; certificates in \texttt{certificates/}).} For every pair \((M,j)\) in
\[\mathcal S=\{(M,1):15\le M\le 10^5\}\cup\{(M,2):44\le M\le10^5\}\cup\{(M,3):87\le M\le10^5\}\]
the scaled Hopf system has a unique zero \((\kappa_j(M),z_j(M))\) in an explicit box. There:
- \(\pm i\omega\) is a simple pair of eigenvalues of \(B\);
- \(\mathrm{Re}\,d\lambda/d\kappa>0\);
- the first Lyapunov coefficient of the system restricted to the mass hyperplane satisfies
\[\boxed{\ \ell_1<0\ \text{ for all }(M,j)\in\mathcal S\setminus\{(15,1)\},\qquad \ell_1>0\ \text{ for }(M,j)=(15,1).\ }\]

Consequences:
- Wherever the remaining hypothesis H2 holds, the Hopf bifurcation at \(\kappa_j(M)\) is \textbf{supercritical}, except
for \(M=15\) where it is \textbf{subcritical}. H2 is certified for \(M\le120\), proved by the paper for
\(M>M_{0,k}\), and observed numerically everywhere in \(\mathcal S\).
- In particular, for every \(M\in[M_{0,1},10^5]\) (and \(M\ge16\)) the first Hopf bifurcation of Theorem 3.3 is
supercritical. Its bifurcating periodic orbits are orbitally asymptotically stable (within the mass hyperplane).

Remarks.
- The first branch exists from \(M=15\), at \(K\approx212.9\), \(\kappa\approx824.5\). For \(M\le14\) no crossing
was found for \(K\le10^5\) (floating-point scan). Branch 2 starts at \(M=44\) (\(K\approx232\)) and branch 3 at
\(M=87\) (\(K\approx265\)). At each onset the crossing comes in from \(K=\infty\).
- \(M=15\) is a genuine exception. \(\ell_1=+1.1952\times10^{-9}\) (unit normalization), certified, with tight
radius. The dense floating-point formula gives the same value. BifurcationKit, an independent code with automatic
differentiation, reports \texttt{SubCritical} there and \texttt{SuperCritical} at \(M=16,17,30\).
- The scale-free value \(\ell_1/M\) (with \(q_1=1\)) is tiny at \(M=15\) (\(+5.8\times10^{-10}\)), against
\(-6.3\times10^{-6}\) at \(M=16\). This suggests that an analytic continuation in \(M\) would have a degenerate
(Bautin) point between 15 and 16; that continuation is not computed here. For the integer problem the first
branch is subcritical exactly at its first value \(M=15\).

Representative certified values. The unit column uses \(q^Hq=1\); the last column is
\(\ell_1/M\) with the normalization \(q_1=1\).

{\def\LTcaptype{none} % do not increment counter
\begin{longtable}[]{@{}
  >{\raggedright\arraybackslash}p{(\linewidth - 10\tabcolsep) * \real{0.1667}}
  >{\raggedright\arraybackslash}p{(\linewidth - 10\tabcolsep) * \real{0.1667}}
  >{\raggedright\arraybackslash}p{(\linewidth - 10\tabcolsep) * \real{0.1667}}
  >{\raggedright\arraybackslash}p{(\linewidth - 10\tabcolsep) * \real{0.1667}}
  >{\raggedright\arraybackslash}p{(\linewidth - 10\tabcolsep) * \real{0.1667}}
  >{\raggedright\arraybackslash}p{(\linewidth - 10\tabcolsep) * \real{0.1667}}@{}}
\toprule\noalign{}
\begin{minipage}[b]{\linewidth}\raggedright
\(M\)
\end{minipage} & \begin{minipage}[b]{\linewidth}\raggedright
\(j\)
\end{minipage} & \begin{minipage}[b]{\linewidth}\raggedright
\(\kappa_j(M)\)
\end{minipage} & \begin{minipage}[b]{\linewidth}\raggedright
\(\omega\)
\end{minipage} & \begin{minipage}[b]{\linewidth}\raggedright
\(\ell_1\) (unit)
\end{minipage} & \begin{minipage}[b]{\linewidth}\raggedright
\(\ell_1/M\) (\(q_1=1\))
\end{minipage} \\
\midrule\noalign{}
\endhead
\bottomrule\noalign{}
\endlastfoot
15 & 1 & 824.53962734 & 87.797 & \(+1.195\times10^{-9}\) & \(+5.8\times10^{-10}\) \\
16 & 1 & 83.095251377 & 8.0421 & \(-1.296\times10^{-5}\) & \(-6.27\times10^{-6}\) \\
30 & 1 & 9.2400641499 & 0.34187 & \(-6.723\times10^{-3}\) & \\
100 & 1 & 3.9348822214 & 2.174e-2 & \(-3.001\times10^{-2}\) & \\
100 & 2 & 15.746601845 & 0.18808 & \(-2.477\times10^{-3}\) & \\
100 & 3 & 121.37739313 & 2.1986 & \(-1.033\times10^{-5}\) & \\
\(10^4\) & 1 & 2.0375686900 & 9.36e-6 & \(-1.542\times10^{-3}\) & \(-0.011015\) \\
\(10^4\) & 2 & 3.7981676429 & & \(-4.525\times10^{-3}\) & \(-0.007971\) \\
\(10^4\) & 3 & 5.5155380850 & & \(-3.798\times10^{-3}\) & \(-0.003637\) \\
\(10^5\) & 1 & 1.9413727095 & 2.778e-7 & \(-4.028\times10^{-4}\) & \(-0.009236\) \\
\(10^5\) & 2 & 3.3593351327 & 1.162e-6 & \(-1.443\times10^{-3}\) & \(-0.007509\) \\
\(10^5\) & 3 & 4.5144663668 & 2.465e-6 & \(-1.347\times10^{-3}\) & \(-0.003525\) \\
\end{longtable}
}

Full data: \texttt{certificates/sweep\_j\{1,2,3\}*.csv.gz} (one line per \(M\)).

\section{\texorpdfstring{The limit \(M\to\infty\)}{The limit M\textbackslash to\textbackslash infty}}

\subsection{7.1 Formal limit problem}

Notation: \(\varepsilon=M^{-1/2}\), \(K=\varepsilon\kappa\), \(x=\ell/M\), \(\tau=\varepsilon^3t\) (paper (35)).
Write \(n_1=1+c\), \(c=\varepsilon\gamma\), \(\gamma=\kappa+\Gamma\).

\begin{enumerate}
\def\labelenumi{\arabic{enumi}.}
\tightlist
\item
  \textbf{Bulk.} Exactly, \(\dot n_\ell=c\,(n_{\ell-1}-n_\ell)+(n_{\ell-1}-2n_\ell+n_{\ell+1})\) for \(2\le\ell\le N\):
  upwind transport with speed \(c\) plus diffusion. With \(n_\ell\approx 1+\varepsilon\kappa+\varepsilon V(x,\tau)\),
  at leading order
  \[V_\tau+\gamma V_x=0 \quad (0<x<1),\qquad V(0,\tau)=\Gamma(\tau).\]
  The inflow value comes from \(n_1=1+\varepsilon\gamma\). Diffusion contributes \(\varepsilon V_{xx}\).
\item
  \textbf{Boundary layer at \(\ell=M\)}, of width \(\ell\)-scale \(1/c=\sqrt M/\gamma\). It is quasi-static: its
  relaxation time is \(O(M)\ll M^{3/2}\). The flux \(J_\ell=cn_\ell-(n_{\ell+1}-n_\ell)\) is constant across it and
  equals \(Kn_M\) at \(\ell=N\). Solving the recurrence gives \(n_\ell=a+C(1+c)^\ell\), where \(a\) is the outer
  value and \(C(1+c)^M=a(\gamma-\kappa)/\kappa\). Its excess mass is
  \(\sum\ell(n_\ell-a)\approx M^2\varepsilon\,a\,\Gamma/(\kappa\gamma)\).
\item
  \textbf{Monomers} relax on the fast time scale \(1/(A(M+2))\). Their equation is replaced by mass conservation,
  \(\sum\ell(n_\ell-A)=0\). With \(a=1+O(\varepsilon)\) this gives
  \[\int_0^1xV\,dx+\frac{\Gamma}{\kappa(\kappa+\Gamma)}=0 .\]
\end{enumerate}

\textbf{Lagrangian clock.} Let \(dT=\gamma\,d\tau\). Along characteristics \(V\) carries its inflow value, so
\(V(x,\cdot)=G(T-x)\) with \(G(T)=\Gamma\). The problem becomes the scalar \emph{renewal equation}
\[\boxed{\ \frac{G(T)}{\kappa(\kappa+G(T))}+\int_0^1 s\,G(T-s)\,ds=0\ }\tag{R}\]
\((R)\) is autonomous, and its only nonlinearity is the boundary-layer term. A monotone time change preserves
orbits, their stability and the side of bifurcation, so criticality can be read on \((R)\).

\textbf{Check against the paper.} With \(G=e^{\zeta T}\), linearizing (R) gives
\(\Delta(\zeta)=\kappa^{-2}+E(\zeta)=0\), where \(E(\zeta)=\int_0^1se^{-\zeta s}ds=(1-e^{-\zeta}(1+\zeta))/\zeta^2\).
Since \(\zeta^2e^\zeta\Delta(\zeta)=Q(\zeta;\kappa)\), this is exactly the paper's limit function \(Q\)
(eq. (87)). Its purely imaginary roots are \(\zeta=\pm it_j\), where \(\tan t_j=t_j\), \(\cos t_j<0\) and
\(\kappa_j^0{}^2=\sqrt{1+t_j^2}-1\).

\subsection{7.2 Landau coefficient of (R) (multiple scales)}

Multiply (R) by \(\kappa^2\):
\[G+\kappa^2LG-G^2/\kappa+G^3/\kappa^2+O(G^4)=0,\qquad (LG)(T)=\int_0^1sG(T-s)ds .\]
Expand with \(\kappa=\kappa_c+\eta^2\tilde\kappa\), \(G=\eta G_1+\eta^2G_2+\eta^3G_3\),
\(G_1=\mathcal A(\eta^2T)e^{i\omega T}+\text{c.c.}\), \(\omega=t_j\), \(D(\zeta):=1+\kappa_c^2E(\zeta)\).
- Order \(\eta\): \(D(i\omega)=0\).
- Order \(\eta^2\): \(G_2=\alpha\mathcal A^2e^{2i\omega T}+\text{c.c.}+\beta|\mathcal A|^2\), with
\(\alpha=1/(\kappa D(2i\omega))\) and \(\beta=2/(\kappa D(0))\).
- Order \(\eta^3\): solvability of the \(e^{i\omega T}\) component, using \(L[\mathcal A(\eta^2T)e^{i\omega T}]
  =e^{i\omega T}(E(i\omega)\mathcal A+\eta^2E'(i\omega)\mathcal A')\) and \(E(i\omega)=-1/\kappa^2\).
Explicitly, using \(\kappa^2=\kappa_c^2+2\kappa_c\eta^2\tilde\kappa+\dots\), the \(e^{i\omega T}\) component of the
order-\(\eta^3\) equation reads
\[\kappa_c^2E'(i\omega)\mathcal A'+2\kappa_c\tilde\kappa E(i\omega)\mathcal A-\frac2{\kappa_c}(\alpha+\beta)|\mathcal A|^2\mathcal A+\frac3{\kappa_c^2}|\mathcal A|^2\mathcal A=0 .\]
Here \((G_1G_2)_{e^{i\omega T}}=(\alpha+\beta)|\mathcal A|^2\mathcal A\) and
\((G_1^3)_{e^{i\omega T}}=3|\mathcal A|^2\mathcal A\).

The result is the amplitude equation
\[\mathcal A'=\sigma\tilde\kappa\mathcal A+\ell\,|\mathcal A|^2\mathcal A,\qquad
\sigma=\frac{2}{\kappa^3E'(i\omega)},\qquad
\ell=\frac{b}{\kappa^4E'(i\omega)},\qquad b=\frac2{D(2i\omega)}+\frac4{D(0)}-3 .\]

\textbf{Exact evaluation at the Hopf points} (\texttt{05\_limit\_coefficient.py}, C1--C3, sympy). Set \(R=\sqrt{1+t^2}=1+\kappa^2\).
The Hopf conditions give \(e^{it}=-(1+it)/R\), and then:
\[E'(it)=\frac1R-i\frac{R+1}{tR(R-1)},\qquad E(2it)=\frac{1-it}{2R^2},\qquad D(0)=\frac{R+1}2,\]
\[\mathrm{Re}\big(b\,\overline{E'(it)}\big)=-\frac{7R^2-34R+23}{R(R+1)(5R-3)} .\]
The last identity is checked exactly in the rational parametrization \(t=(s^2-1)/(2s)\), \(R=(s^2+1)/(2s)\).

\textbf{Signs.}
- \(\mathrm{Re}\,\sigma>0\), since \(\mathrm{Re}\,E'(it)=1/R>0\). This is the transversal crossing, consistent with
paper Lemma 7.3.
- \(7R^2-34R+23>0\) iff \(R>(17+8\sqrt2)/7\approx4.0448\). Every Hopf point has \(t_j\ge t_1>3.92\) (interval
check: \(\tan t-t\) is increasing on \((\pi,3\pi/2)\) and negative at \(3.92\)), hence \(R_j>4.0455\).
- Therefore \(\mathrm{Re}\,\ell_j<0\) \textbf{for every \(j\)}. Equivalently, the sign is that of
\(-(7\kappa^4-20\kappa^2-4)\), negative since \(\kappa_j^0\ge\kappa_{\rm cr}=1.8982>1.7449\).

\textbf{Proposition B.} Every Hopf bifurcation of the limit equation (R), at \(\kappa=\kappa_j^0\), \(j\ge1\), has a
negative Landau coefficient. Hence it is supercritical in the multiple-scales sense.
- The algebra is exact.
- The multiple-scales coefficient is the standard first Lyapunov coefficient, up to a positive factor; the rigorous
Hopf theory for renewal equations is in Diekmann, Getto and Gyllenberg (2008).
- A direct simulation of (R) confirms the amplitude law. At \(\kappa=\kappa_1^0+\tilde\kappa\), a stable cycle forms
with amplitude ratio (simulation/prediction) \(1.038\) at \(\tilde\kappa=0.004\) and \(1.018\) at \(0.002\)
(claim C4b).

\textbf{Prediction for the discrete \(\ell_1\)} (heuristic matching of normalizations; confirmed numerically in §8).
Normalize \(q_1=1\), so \(v_1=2\,\mathrm{Re}\,z_{\rm Kuz}\). Use \(v_1=\varepsilon G\), \(T\approx\kappa\tau\) and
\(\tau=\varepsilon^3t\). Then \(c_1=\varepsilon\kappa\ell\) and \(\omega=\varepsilon^3\kappa t_j\), so
\[\frac{\ell_1(q_1=1)}{M}\ \longrightarrow\ L_j:=\frac{\mathrm{Re}\,\ell_j}{t_j}
=-\frac{R(R-1)(7R^2-34R+23)}{t\,(R+1)(5R-3)\big((R-1)^3+R+1\big)}\Big|_{t=t_j}.\]
Values: \(L_1=-0.0093114\), \(L_2=-0.0079410\), \(L_3=-0.0037777\), \(L_4=-0.0021661\), \(L_5=-0.0013977\).
Also \(L_j\sim-7/(5t_j^2)\to0^-\).

\section{Convergence evidence (Observation C)}

At sample points \(M=10^4,3\cdot10^4,\dots,10^8\) and \(j=1,2,3\), the Hopf point and \(\ell_1\) are certified as in
§5 (\texttt{06\_large\_M\_samples.py}):
- \(\ell_1<0\) at every sample point;
- the relative deviation of \(\ell_1(q_1=1)/M\) from \(L_j\) at \(M=10^8\) is \(1.3\times10^{-3}\) (\(j=1\)),
\(3.6\times10^{-3}\) (\(j=2\)) and \(6.7\times10^{-3}\) (\(j=3\));
- the deviation decays like \(M^{-0.44}\), consistent with an \(O(M^{-1/2})=O(\varepsilon)\) correction.

So the discrete coefficient converges to the limit value predicted by §7, which independently confirms both the
formal limit and the closed form.

\section{\texorpdfstring{Validation of the machinery (\texttt{02\_validate\_l1.py})}{Validation of the machinery (02\_validate\_l1.py)}}

\begin{itemize}
\tightlist
\item
  \textbf{V1.} The generic Kuznetsov routine reproduces the Guckenheimer--Holmes coefficient on a planar test system to
  \(10^{-12}\).
\item
  \textbf{V2.} Direct stiff simulation of the full system at \(M=30\), \(K=K_c(1+\delta)\). A small perturbation converges
  to a small stable cycle. The \(n_1\) half-amplitude divided by the normal-form prediction
  \(2|q_1|\sqrt{-\mu/(\omega\ell_1)}\) is \(1.0096,1.0047,1.0024,1.0012\) for \(\delta=4,2,1,0.5\%\): the ratio tends
  to 1 linearly in \(\delta\).
\item
  \textbf{V3.} BifurcationKit 0.8 (Julia) classifies the Hopf points as follows: \(M=15\) SubCritical (\(\mathrm{Re}\,b=
  +1.22\times10^{-7}\)); \(M=16\), \(17\) and \(30\) SuperCritical. Its \(K_c\) agree with ours to \(10^{-9}\). This
  code is independent: reduced coordinates, automatic differentiation, its own normal-form code.
\item
  \textbf{V4} (\texttt{bk\_po\_branch.jl}, log in \texttt{certificates/logs/bk\_po\_branch.log}). The branch of periodic orbits from the
  Hopf point is continued by orthogonal collocation (BifurcationKit).

  \begin{itemize}
  \tightlist
  \item
    \(M=16\): the branch leaves towards \(K>K_c\) with amplitude\(^2\propto K-K_c\): supercritical.
  \item
    \(M=15\): the branch leaves towards \(K<K_c\) with amplitude\(^2\propto K_c-K\) (ratios 12.4 vs 12.8):
    subcritical. It then folds back towards \(K>K_c\) at \(n_M\)-amplitude \(\approx1.2\), a saddle-node of cycles,
    consistent with the tiny \(\ell_1\).
  \end{itemize}
\end{itemize}

\section{What is proved, what is not, and points for the verifier}

\textbf{Proved (computer-assisted, ball arithmetic).}
- Theorem A: the sign of \(\ell_1\) for every \((M,j)\in\mathcal S\), with simplicity and transversality.
- H2 and the branch label for \(M\le120\).

\textbf{Exact (symbolic).}
- The Landau coefficient of the limit equation (R) and its sign for every \(j\) (Proposition B, algebraic part).

\textbf{Formal or not proved.}
1. \emph{Large \(M\) is not covered rigorously.} There is no proof that \(\ell_1(q_1=1)/M\to L_j\); §8 is numerical
evidence plus certified samples. A complete answer for all \(M>M_{0,k}\) needs this convergence with an explicit
rate, together with Theorem A up to the resulting threshold. An \(\varepsilon\)-graded interval evaluation of the
closed form is not directly possible: the inner products contain an exact cancellation of relative order
\(\varepsilon^{-3}\) between the row-1 term and the bulk sum (mass conservation). A proof would have to reorganize
the computation around the continuum structure of §7.
2. \emph{Derivation of (R)} (§7.1): formal (quasi-static boundary layer, slaved monomers, neglected diffusion). It
reproduces the paper's \(Q\) at linear order. The match of §8 (to \(10^{-3}\) at \(M=10^8\)) supports it at
cubic order.
3. \emph{Multiple-scales coefficient as the Lyapunov coefficient of (R):} standard but not re-proved here (renewal
equations: sun-star calculus).
4. \emph{H2 for \(120<M\)} and the branch label \(j\) for \(M>300\): floating-point evidence only, plus continuation.
The paper covers H2 for \(M>M_{0,k}\).
5. \emph{Onset values} \(M=15,44,87\): floating-point scans with \(K\le10^5\). This affects only the \emph{existence} of the
branch below these values, not the certified signs.
6. \(j\ge4\) is not certified. Floating point gives \(\ell_1<0\) for \(j\le6\) at \(M=150\)--\(300\), and the limit
coefficient is negative for every \(j\).

\textbf{Points for an independent verifier.}
- (a) Re-derive (2.1) and the role of \(H\) in (3.1): the choice \(w_{11}\in H\).
- (b) Check the ansatz argument in §4, in particular the completeness of the family and row 1 versus the mass
condition.
- (c) Re-run \texttt{04\_verify\_certificates.py\ 1\ ...}. It re-proves every line from the stored midpoints alone.
- (d) Re-derive the boundary-layer mass \(\Gamma/(\kappa(\kappa+\Gamma))\) and the Lagrangian reduction to (R).
- (e) Check the multiple-scales algebra (signs of \(c_3\)), using \texttt{05\_limit\_coefficient.py}, claims C2 and C4b.
- (f) Confirm \(M=15\) independently. Already done here by BifurcationKit's normal form (V3) and its periodic-orbit
continuation (V4); a third, independent code would be welcome.

\section{Reproduction}

From the project root in Git Bash: \texttt{outils/wsl\ "cd\ chemistry-open-problems/problems/becker-doring-atomization-hopf/scripts\ \&\&\ python\ \textless{}script\textgreater{}"}.

{\def\LTcaptype{none} % do not increment counter
\begin{longtable}[]{@{}
  >{\raggedright\arraybackslash}p{(\linewidth - 4\tabcolsep) * \real{0.3333}}
  >{\raggedright\arraybackslash}p{(\linewidth - 4\tabcolsep) * \real{0.3333}}
  >{\raggedright\arraybackslash}p{(\linewidth - 4\tabcolsep) * \real{0.3333}}@{}}
\toprule\noalign{}
\begin{minipage}[b]{\linewidth}\raggedright
script
\end{minipage} & \begin{minipage}[b]{\linewidth}\raggedright
content
\end{minipage} & \begin{minipage}[b]{\linewidth}\raggedright
runtime
\end{minipage} \\
\midrule\noalign{}
\endhead
\bottomrule\noalign{}
\endlastfoot
\texttt{01\_check\_model.py} & model identities, paper's eigenvalues and Table 1, closed form = dense & 1 min \\
\texttt{02\_validate\_l1.py} & planar test, stiff simulation (M = 30), BifurcationKit cross-check & 10 min \\
\texttt{03\_certify\_range.py\ J\ MMIN\ MMAX\ OUT.csv\ {[}SEED.csv{]}} & certified sweep (Theorem A) & 12--30 ms per \(M\) \\
\texttt{04\_verify\_certificates.py\ STRIDE\ files...} & independent re-check of the sweep files & 20 ms per line \\
\texttt{05\_limit\_coefficient.py} & limit problem, exact algebra, simulation of (R) & 1 min \\
\texttt{06\_large\_M\_samples.py} & certified samples up to \(M=10^8\), convergence to \(L_j\) & 2 min \\
\texttt{07\_spectrum\_H2.py\ MMAX\ files...} & rigorous spectra, H2 and branch label for \(M\le\) MMAX & 15 min \\
\texttt{bk\_check.jl\ M\ Kc\ ...} & BifurcationKit normal form (independent) & 2 min \\
\texttt{bk\_po\_branch.jl} & BifurcationKit periodic-orbit branch at \(M=15,16\) (direction of bifurcation) & 3 min \\
\end{longtable}
}

\end{document}
